% GENERATED FILE. Edit canonical lessons, then run npm run generate:books.
% canonical-source-hash: fnv1a64:13e26851245e1126
% canonical-lessons: LA-C435-navigium, LA-C435-navicula, LA-C435-scapha, LA-C435-linter, LA-C435-ratis

\chapter{Vessel, Little boat, Skiff, Punt, Raft}
\label{ch:la-vessel-little-boat-skiff-punt-raft}
\hlchaptermodality{435}

\begin{chapteropening}
\textbf{By the end of this chapter:} \emph{I can name boats of every size: a vessel, a little boat, a skiff, a punt and a raft.}

Complete the last lesson of chapter 435: I can name boats of every size: a vessel, a little boat, a skiff, a punt and a raft.
\end{chapteropening}

\section[nāvigium, nāvigiī]{nāvigium, nāvigiī — a vessel, a craft (any ship or boat)}

\label{lesson:LA-C435-navigium}



\begin{quote}
Before the new one: say the Latin for to neigh, then the Latin for a muleteer.
\end{quote}

\subsection*{You'll want to know: nāvigium, nāvigiī}
\textbf{nāvigium, nāvigiī} — \textquotedblleft{}a vessel, a craft (any ship or boat)\textquotedblright{}.

\begin{quote}
In portū multa nāvigia sunt, magna et parva.
\end{quote}

\begin{quote}
\emph{In the harbour there are many vessels, large and small.}
\end{quote}

Second declension, neuter: \emph{nāvigium, nāvigiī n.} A general word for anything that sails.

\begin{cousinweb}
\emph{nāvigium} is built on \emph{nāvigō} 'to sail' (EtymDB; Lewis \& Short), which comes from \emph{nāvis} 'ship', inherited from PIE \emph{néh\textsubscript{2}us} 'boat'.
\end{cousinweb}

\begin{grammarlens}[title={What you've built}]
One more word for telling, explaining and describing at B1.
\end{grammarlens}

\subsection*{Your turn}
\begin{itemize}
\raggedright
  \item \emph{Say it:} \emph{nāvigium, nāvigiī}
  \item \emph{Say it:} \emph{nāvigium, nāvigiī}, once more
  \item \emph{Say it:} a sentence of your own with \emph{nāvigium, nāvigiī}
  \item \emph{Recall:} say the Latin for to neigh, then the Latin for a muleteer, then say \emph{nāvigium, nāvigiī} again
\end{itemize}

\begin{tcolorbox}[breakable,colback=teal!4,colframe=teal!35!black,title={Before you move on}]
What does nāvigium, nāvigiī mean? (\textquotedblleft{}A vessel, a craft (any ship or boat)\textquotedblright{}.) Say it once more.
\end{tcolorbox}
\section[nāvicula, nāviculae]{nāvicula, nāviculae — a little boat, a small ship}

\label{lesson:LA-C435-navicula}



\begin{quote}
Before the new one: say the Latin for a muleteer, then the Latin for a vessel.
\end{quote}

\subsection*{You'll want to know: nāvicula, nāviculae}
\textbf{nāvicula, nāviculae} — \textquotedblleft{}a little boat, a small ship\textquotedblright{}.

\begin{quote}
Piscātor nāviculā in flūmine piscēs capit.
\end{quote}

\begin{quote}
\emph{The fisherman catches fish from a little boat on the river.}
\end{quote}

First declension, feminine: \emph{nāvicula, nāviculae f.}, a diminutive of \emph{nāvis}. English \emph{navicular} 'boat-shaped' comes from it.

\begin{cousinweb}
\emph{nāvicula} is a diminutive of \emph{nāvis} 'ship' (Lewis \& Short), made with the ending \emph{-cula}. \emph{nāvis} is inherited from PIE \emph{néh\textsubscript{2}us} 'boat' (EtymDB).
\end{cousinweb}

\begin{grammarlens}[title={What you've built}]
One more word for telling, explaining and describing at B1.
\end{grammarlens}

\subsection*{Your turn}
\begin{itemize}
\raggedright
  \item \emph{Say it:} \emph{nāvicula, nāviculae}
  \item \emph{Say it:} \emph{nāvicula, nāviculae}, once more
  \item \emph{Say it:} a sentence of your own with \emph{nāvicula, nāviculae}
  \item \emph{Recall:} say the Latin for a muleteer, then the Latin for a vessel, then say \emph{nāvicula, nāviculae} again
\end{itemize}

\begin{tcolorbox}[breakable,colback=teal!4,colframe=teal!35!black,title={Before you move on}]
What does nāvicula, nāviculae mean? (\textquotedblleft{}A little boat, a small ship\textquotedblright{}.) Say it once more.
\end{tcolorbox}
\section[scapha, scaphae]{scapha, scaphae — a skiff, a ship's boat}

\label{lesson:LA-C435-scapha}



\begin{quote}
Before the new one: say the Latin for a vessel, then the Latin for a little boat.
\end{quote}

\subsection*{You'll want to know: scapha, scaphae}
\textbf{scapha, scaphae} — \textquotedblleft{}a skiff, a ship's boat\textquotedblright{}.

\begin{quote}
Nautae ē nāve in scapham dēscendunt et ad lītus remīs eunt.
\end{quote}

\begin{quote}
\emph{The sailors climb down from the ship into the skiff and row to the shore.}
\end{quote}

First declension, feminine: \emph{scapha, scaphae f.} The small boat a ship carried or towed, for going ashore.

\begin{cousinweb}
\emph{scapha} was borrowed from Greek \emph{skaphē} 'light boat, skiff' (EtymDB); Lewis \& Short compare \emph{scabō} 'to scratch, to scrape', as of a hollowed-out boat.
\end{cousinweb}

\begin{grammarlens}[title={What you've built}]
One more word for telling, explaining and describing at B1.
\end{grammarlens}

\subsection*{Your turn}
\begin{itemize}
\raggedright
  \item \emph{Say it:} \emph{scapha, scaphae}
  \item \emph{Say it:} \emph{scapha, scaphae}, once more
  \item \emph{Say it:} a sentence of your own with \emph{scapha, scaphae}
  \item \emph{Recall:} say the Latin for a vessel, then the Latin for a little boat, then say \emph{scapha, scaphae} again
\end{itemize}

\begin{tcolorbox}[breakable,colback=teal!4,colframe=teal!35!black,title={Before you move on}]
What does scapha, scaphae mean? (\textquotedblleft{}A skiff, a ship's boat\textquotedblright{}.) Say it once more.
\end{tcolorbox}
\section[linter, lintris]{linter, lintris — a punt, a small flat boat; a trough}

\label{lesson:LA-C435-linter}



\begin{quote}
Before the new one: say the Latin for a little boat, then the Latin for a skiff.
\end{quote}

\subsection*{You'll want to know: linter, lintris}
\textbf{linter, lintris} — \textquotedblleft{}a punt, a small flat boat; a trough\textquotedblright{}.

\begin{quote}
Agricolae lintribus flūmen trānseunt.
\end{quote}

\begin{quote}
\emph{The farmers cross the river in punts.}
\end{quote}

Third declension, feminine: \emph{linter, lintris f.} A flat river boat; also a wooden tub or trough.

\begin{cousinweb}
The origin of \emph{linter} is not settled: Lewis \& Short connect it, doubtfully, with Greek \emph{plynō} 'to wash', a washtub before a boat, and EtymDB lists it without a parent.
\end{cousinweb}

\begin{grammarlens}[title={What you've built}]
One more word for telling, explaining and describing at B1.
\end{grammarlens}

\subsection*{Your turn}
\begin{itemize}
\raggedright
  \item \emph{Say it:} \emph{linter, lintris}
  \item \emph{Say it:} \emph{linter, lintris}, once more
  \item \emph{Say it:} a sentence of your own with \emph{linter, lintris}
  \item \emph{Recall:} say the Latin for a little boat, then the Latin for a skiff, then say \emph{linter, lintris} again
\end{itemize}

\begin{tcolorbox}[breakable,colback=teal!4,colframe=teal!35!black,title={Before you move on}]
What does linter, lintris mean? (\textquotedblleft{}A punt, a small flat boat; a trough\textquotedblright{}.) Say it once more.
\end{tcolorbox}
\section[ratis, ratis]{ratis, ratis — a raft; (in poetry) a boat}

\label{lesson:LA-C435-ratis}



\begin{quote}
Before the new one: say the Latin for a skiff, then the Latin for a punt.
\end{quote}

\subsection*{You'll want to know: ratis, ratis}
\textbf{ratis, ratis} — \textquotedblleft{}a raft; (in poetry) a boat\textquotedblright{}.

\begin{quote}
Mīlitēs ratēs ex arboribus faciunt et flūmen trānseunt.
\end{quote}

\begin{quote}
\emph{The soldiers make rafts out of trees and cross the river.}
\end{quote}

Third declension, feminine: \emph{ratis, ratis f.} Logs tied together; poets use it for any ship.

\begin{cousinweb}
\emph{ratis} is inherited. Lewis \& Short connect it with \emph{rēmus} 'oar' and Greek \emph{eretēs} 'rower', but EtymDB gives no parent, so the link is hedged.
\end{cousinweb}

\begin{grammarlens}[title={What you've built}]
One more word for telling, explaining and describing at B1.
\end{grammarlens}

\subsection*{Your turn}
\begin{itemize}
\raggedright
  \item \emph{Say it:} \emph{ratis, ratis}
  \item \emph{Say it:} \emph{ratis, ratis}, once more
  \item \emph{Say it:} a sentence of your own with \emph{ratis, ratis}
  \item \emph{Recall:} say the Latin for a skiff, then the Latin for a punt, then say \emph{ratis, ratis} again
\end{itemize}

\begin{tcolorbox}[breakable,colback=teal!4,colframe=teal!35!black,title={Before you move on}]
What does ratis, ratis mean? (\textquotedblleft{}A raft; (in poetry) a boat\textquotedblright{}.) Say it once more.
\end{tcolorbox}
